% Banco canónico de veinte pruebas. Las lecturas son soluciones del autor,
% no resultados de una prueba independiente de lectura.
\KalaCaso{T01}{Matemáticas}
 {Para cada natural hay un natural mayor.}
 {\KalaUniversal{x}{\mathbb{N}}{\KalaExistencia{y}{\mathbb{N}}{\KalaMat{x<y}}}}
 {\(\mathbb N=\{0,1,2,\ldots\}\). Las variables quedan vinculadas por sus cuantificadores.}
 {Para todo \(x\) natural existe \(y\) natural con \(x<y\); \(y\) puede depender de \(x\).}
 {Mover el existencial delante del universal exigiría un natural mayor que todos. La expresión actual es verdadera: sirve \(y=x+1\).}
 {control de compatibilidad con E01.}

\KalaCaso{T02}{Lógica y filosofía}
 {Negar que dos proposiciones sean ambas verdaderas.}
 {\KalaNegacion{\KalaConjuncion{\KalaMat{P}}{\KalaMat{Q}}}}
 {\(P,Q\) son proposiciones arbitrarias.}
 {No es verdad que \(P\) y \(Q\) sean ambas verdaderas.}
 {Con \(P=\text{verdadero}\) y \(Q=\text{falso}\), T02 es verdadera y T03 es falsa.}
 {negación con alcance sobre una combinación.}

\KalaCaso{T03}{Lógica y filosofía}
 {Afirmar que P es falsa y Q verdadera.}
 {\KalaConjuncion{\KalaNegacion{\KalaMat{P}}}{\KalaMat{Q}}}
 {Los mismos significados de \(P,Q\) que en T02.}
 {Se cumple la negación de \(P\) y se cumple \(Q\).}
 {El grupo de negación contiene solo \(P\). T02 y T03 no son equivalentes.}
 {comparación de alcances.}

\KalaCaso{T04}{Lectura y vida cotidiana}
 {A cada persona le gusta al menos un libro.}
 {\KalaTodo{x}{\KalaConcepto{1}}{\KalaExiste{y}{\KalaConcepto{2}}{\KalaGusta{\KalaVariable{x}}{\KalaVariable{y}}}}}
 {C1: personas del estudio; C2: libros del estudio. R7(persona, libro): a la persona le gusta ese libro según el criterio declarado.}
 {Cada persona puede elegir un libro diferente. No se exige un favorito único.}
 {Dos personas y dos libros, con gustos separados, hacen verdadera T04 y falsa T05.}
 {Extensión: conceptos indexados, dominios no numéricos y R7; cuantificadores conservados.}

\KalaCaso{T05}{Lectura y vida cotidiana}
 {Hay un mismo libro que les gusta a todas las personas.}
 {\KalaExiste{y}{\KalaConcepto{2}}{\KalaTodo{x}{\KalaConcepto{1}}{\KalaGusta{\KalaVariable{x}}{\KalaVariable{y}}}}}
 {Los mismos C1, C2 y R7 que en T04. En la comparación se usan dos personas y dos libros.}
 {Existe un libro del estudio que satisface la relación de gusto con cada persona.}
 {El testigo \(y\) se elige antes de recorrer las personas. No puede cambiar con \(x\).}
 {Extensión léxica; prueba de orden de cuantificadores.}

\KalaCaso{T06}{Arquitectura de computadoras}
 {El procesador lee el dato D de la dirección A.}
 {\KalaRelacion{2}{\KalaReferencia{1}\KalaSeparador\KalaReferencia{2}\KalaSeparador\KalaReferencia{3}}}
 {Referentes: 1 procesador; 2 dirección A; 3 dato D. Un único evento abstracto de lectura.}
 {El agente 1 lee el dato 3 desde la dirección 2.}
 {Intercambiar los argumentos 2 y 3 cambia sus papeles. No se describe una escritura ni un modelo de caché.}
 {relación ternaria R2.}

\KalaCaso{T07}{Biología}
 {La estructura A es parte propia de la célula B.}
 {\KalaRelacion{3}{\KalaReferencia{1}\KalaSeparador\KalaReferencia{2}}}
 {Referentes: 1 estructura A; 2 célula B, objetos concretos de un modelo ficticio. R3 es parte propia.}
 {A forma parte de B y A no es idéntica a B.}
 {No afirma que A sea un elemento de un conjunto B; esa es la relación de T08.}
 {distinción entre parte y miembro.}

\KalaCaso{T08}{Colecciones y matemáticas discretas}
 {El objeto A pertenece al conjunto formado por A y B.}
 {\KalaMiembro{\KalaReferencia{1}}{\KalaColeccion{1}{\KalaReferencia{1}\KalaSeparador\KalaReferencia{2}}}}
 {Referentes: 1 objeto A; 2 objeto B, distintos en este ejemplo. El índice 1 de K06 indica conjunto; el índice 2 indicaría lista.}
 {A es un miembro de \(\{A,B\}\).}
 {El conjunto \(\{A,B\}\) coincide con \(\{B,A\}\). Las listas \([A,B]\) y \([B,A]\) son diferentes. R5 solo acepta conjuntos o extensiones de conceptos.}
 {Extensión: construcción de colecciones y pertenencia R5.}

\KalaCaso{T09}{Historia}
 {El evento A terminó antes de comenzar B.}
 {\KalaRelacion{1}{\KalaReferencia{1}\KalaSeparador\KalaReferencia{2}}}
 {Referentes: 1 evento A, intervalo \([1,2]\); 2 evento B, intervalo \([3,4]\), misma escala temporal ficticia.}
 {El final de A precede estrictamente al inicio de B.}
 {La cronología no basta para obtener una afirmación causal como T10.}
 {relación temporal R1.}

\KalaCaso{T10}{Modelos causales}
 {En el modelo M, intervenir sobre X puede cambiar Y.}
 {\KalaRelacion{6}{\KalaReferencia{1}\KalaSeparador\KalaReferencia{2}\KalaSeparador\KalaReferencia{3}}}
 {Referentes: 1 variable X; 2 variable Y; 3 modelo determinista acíclico M con \(X:=U\), \(Y:=X\), \(U\in\{0,1\}\). Se compara el mismo contexto U bajo \(\operatorname{do}(X=0)\) y \(\operatorname{do}(X=1)\).}
 {En ese modelo, existe un contexto y dos intervenciones sobre X que producen valores distintos de Y.}
 {No afirma que una correlación observada demuestre causalidad, ni formula una teoría general de causa de eventos.}
 {Extensión: R6, con modelo y significado de intervención explícitos.}

\KalaCaso{T11}{Física experimental}
 {Atribuir una medición y su incertidumbre estándar a un registro.}
 {\KalaFuente{\KalaReferencia{2}}{\KalaMat{m(S_1)=5{,}0\,\mathrm{g};\quad u(m)=0{,}1\,\mathrm{g}}}}
 {Referente 2: registro ficticio de la muestra \(S_1\), con procedimiento y condiciones identificados. \(u(m)\) denota incertidumbre estándar.}
 {El registro 2 informa esos valores.}
 {K16 no certifica su verdad. La incertidumbre indicada no es un error conocido ni un intervalo de confianza.}
 {atribución y notación científica convencional.}

\KalaCaso{T12}{Alternativas lógicas}
 {Afirmar que P o Q se cumple, permitiendo ambas.}
 {\KalaAlternativa{\KalaMat{P}}{\KalaMat{Q}}}
 {\(P,Q\) son proposiciones; «o» se usa aquí en sentido inclusivo.}
 {Al menos una es verdadera.}
 {Si ambas son verdaderas, T12 sigue siendo verdadera. La exclusión se escribe combinando T12 con la negación de su conjunción.}
 {Extensión: K17, disyunción inclusiva.}

\KalaCaso{T13}{Posibilidad y modelos}
 {P es posible según el modelo modal 1.}
 {\KalaPosible{1}{\KalaMat{P}}}
 {Modelo 1: estados \(w_0,w_1\); estado actual \(w_0\); accesibilidad \(A_1=\{(w_0,w_1)\}\). P es falsa en \(w_0\) y verdadera en \(w_1\).}
 {Hay un estado accesible desde el actual donde P es verdadera.}
 {T13 es verdadera aunque P sea falsa en el estado actual. No asigna una probabilidad.}
 {Extensión: K18. Se limita al fragmento proposicional con modelo declarado.}

\KalaCaso{T14}{Preguntas y diálogo escrito}
 {Preguntar si el evento A precedió al B.}
 {\KalaPregunta{\KalaRelacion{1}{\KalaReferencia{1}\KalaSeparador\KalaReferencia{2}}}}
 {Referentes: 1 evento A; 2 evento B, con calendario común. En este ejercicio no se conocen sus fechas.}
 {¿Terminó A antes de comenzar B?}
 {Se solicita una respuesta; no se afirma la relación. «No sé» es un estado de información, distinto de negar R1.}
 {Extensión: K19. Pregunta polar en el nivel superior.}

\KalaCaso{T15}{Diseño de algoritmos}
 {Inicializar x en cero y aumentarlo mientras sea menor que tres.}
 {\KalaSecuencia{\KalaAsignar{x}{0}}{\KalaMientras{\KalaMat{x<3}}{\KalaAsignar{x}{x+1}}}}
 {Programa determinista sobre enteros matemáticos sin desbordamiento. La condición se comprueba antes del cuerpo.}
 {La traza de x después de inicializar es \(0,1,2,3\); después se detiene.}
 {Un bucle mientras con condición falsa ejecuta cero iteraciones. Una repetición hasta comprueba después de una primera ejecución. K13, el condicional lógico, no ejecuta pasos.}
 {Extensión: asignación indexada K04, secuencia K21 y mientras K22.}

\KalaCaso{T16}{Ingeniería y composición de conceptos}
 {Construir el concepto «sensores rojos y activos».}
 {\KalaSeleccion{x}{\KalaConcepto{3}}{\KalaConjuncion{\KalaPosee{\KalaVariable{x}}{1}}{\KalaPosee{\KalaVariable{x}}{2}}}}
 {C3: sensores del estudio. Propiedad 1: rojo según clasificación declarada. Propiedad 2: activo en el instante de inspección \(t_0\), común a toda la expresión.}
 {El conjunto de sensores que satisfacen ambas propiedades en el contexto indicado.}
 {Construye un concepto; no afirma que existan sensores que lo satisfagan. Para afirmar existencia se cuantifica sobre este resultado.}
 {Extensión: K20, propiedades indexadas y R8. Se reutilizan los conectivos.}

\KalaCaso{T17}{Química y notación especializada}
 {Registrar la proporción estequiométrica de formación de agua a partir de hidrógeno y oxígeno.}
 {\KalaMat{2\,\mathrm{H}_2+\mathrm{O}_2\longrightarrow 2\,\mathrm{H}_2\mathrm{O}}}
 {Región convencional de química. Los coeficientes indican proporciones estequiométricas; se omiten estados físicos y condiciones.}
 {Dos unidades estequiométricas de hidrógeno molecular y una de oxígeno molecular corresponden a dos de agua.}
 {La ecuación conserva cuatro átomos de H y dos de O a cada lado. No describe velocidad, mecanismo ni condiciones de realización.}
 {Interoperabilidad mediante región convencional; no se crea otro alfabeto químico.}

\KalaCaso{T18}{Geografía}
 {A está al norte de B en el marco geográfico F.}
 {\KalaRelacion{9}{\KalaReferencia{1}\KalaSeparador\KalaReferencia{2}\KalaSeparador\KalaReferencia{3}}}
 {Referentes: 1 lugar A con latitud \(12^\circ\); 2 lugar B con latitud \(10^\circ\); 3 marco F, mismo datum y época, latitud positiva hacia el norte. Lugares ficticios.}
 {La latitud de A es mayor que la de B en F.}
 {No afirma una distancia ni compara longitudes. Cambiar de marco exige declarar otro referente.}
 {Extensión: R9 con marco explícito.}

\KalaCaso{T19}{Economía descriptiva}
 {Atribuir a una tabla un aumento nominal del 20 por ciento.}
 {\KalaFuente{\KalaReferencia{1}}{\KalaMat{p_0=10\,\mathrm{UM};\quad p_1=12\,\mathrm{UM};\quad (p_1-p_0)/p_0=0{,}20}}}
 {Referente 1: tabla ficticia del precio del mismo bien, misma unidad de cantidad, en dos periodos. UM es una misma unidad monetaria; no hay ajuste por inflación.}
 {Según la tabla, el precio nominal pasa de 10 a 12 UM, un aumento relativo de 0,20.}
 {Un cambio nominal no establece un cambio real ni su causa. Es un ejemplo de representación de datos ficticios.}
 {fuente, cantidades y comparación convencional.}

\KalaCaso{T20}{Música y acústica}
 {Construir el concepto «tonos con frecuencia fundamental mayor que 440 Hz».}
 {\KalaSeleccion{x}{\KalaConcepto{5}}{\KalaMat{f(x)>440\,\mathrm{Hz}}}}
 {C5: tonos estables del estudio. \(f\) asigna a cada tono su frecuencia fundamental durante el mismo intervalo de observación.}
 {El conjunto de tonos del dominio cuya frecuencia fundamental supera 440 Hz.}
 {El umbral es estricto; 440 Hz queda fuera. No se selecciona por volumen. Un dominio vacío produce un conjunto vacío.}
 {Extensión: K20 con predicado convencional; reutilización de la misma composición en otro campo.}
